\section{Introduction}

% === ROUND 2 --- §1 goals & contents (Round 3 writes the text to this) ====
% Section goal: in ~1.5 pages the reader can restate the PROBLEM, the HYPOTHESIS,
% and the CONTRIBUTIONS, and can picture the setup. Audience = ML readers, not
% brain experts: gloss neuro terms on first use, assume ML terms.
% Beat order + per-beat job:
%  1. Instrument + inversion [Fig. pipeline] -- GOAL: what brain alignment is, and
%     the move from MEASURING it to USING it. CONTAINS: the linear-map instrument;
%     that the field only measures; the inversion (alignment as a training term).
%  2. Problem (boxed) -- GOAL: reader states the problem in one sentence. CONTAINS:
%     distillation discards the middle-layer structure alignment lives in; can an
%     alignment-guided objective preserve/improve it at a matched student budget,
%     and does it buy anything downstream; why it matters (the open 2025-26 case).
%  3. Hypothesis (labelled) -- GOAL: one falsifiable hypothesis. CONTAINS: the brain
%     term is a weak prior / regulariser, not a second teacher; prediction = helps
%     where the text signal is weak (small/compressed/scarce), washes out with data.
%     One line of intuition; the math lives in §3.3.
%  4. Positioning [Fig. 2x2] -- GOAL: locate the contribution as the one empty cell.
%     CONTAINS: the two axes; three filled cells (one cite each); the empty cell.
%  5. Contributions -- GOAL: a scannable explicit list. CONTAINS: the six items
%     (problem+positioning; weak-prior theory; artifact-resistant protocol; reality
%     result; no-preservation result; weak quality-entangled lever). No Q3--Q4 outcomes.
%  6. Program + established -- GOAL: the dependent-question roadmap and what is
%     settled so far. CONTAINS: the question sequence; the established results stated
%     QUALITATIVELY (the numbers live in §4, not here).
% The text below is the v0.2 draft; Round 3 reworks it into this order + the doctrine.
% ===========================================================================

Brain alignment has become a way to study what language models and the human language system share \autocite{tuckute2024,oota2023,merlin2024}.
Linear encoding studies ask how well a model's internal activations predict human brain responses to the same text, and the resulting score has been used to compare architectures, layers, training objectives, scale, and compression \autocite{oota2023,aw2023,gao2024,alkhamissi2025,oota2026}.
Recent work also shows that the score is fragile: it can reflect linguistic structure, but it can also be inflated by low-level stimulus features, shallow acoustic cues, or weak evaluation splits \autocite{oota2024,hadidi2026,jia2026}.
Together, this literature has made brain alignment useful as a measurement.
This thesis asks whether the same signal can be useful as a training signal.
In particular, we study model compression: during knowledge distillation, can a small student be guided by a brain-alignment objective, and does that guidance preserve anything ordinary text-only compression loses?

\begin{figure}[t]
  \centering
  \resizebox{0.98\textwidth}{!}{\begin{tikzpicture}[
  font=\small,
  >=Latex,
  node distance=0.85cm and 0.9cm,
  stage/.style={draw=black!55, rounded corners=3pt, fill=gray!6, minimum height=1.05cm, align=center, inner sep=6pt},
  object/.style={draw=black!55, rounded corners=3pt, fill=white, minimum height=0.78cm, align=center, inner sep=5pt},
  signal/.style={draw=black!45, rounded corners=2pt, fill=blue!6, minimum height=0.66cm, align=center, inner sep=4pt},
  loss/.style={draw=black!50, rounded corners=2pt, fill=orange!10, minimum height=0.66cm, align=center, inner sep=4pt},
  arrow/.style={->, line width=0.55pt, draw=black!70},
  softarrow/.style={->, line width=0.5pt, draw=black!45},
]
  \node[stage, text width=2.55cm] (stim) at (-5.0,0) {\textbf{Shared text}\\same story or sentence};
  \node[object, text width=3.05cm] (human) at (-0.55,1.05) {brain response\\to same text\\fMRI $Y$};
  \node[object, text width=3.05cm] (lm) at (-0.55,-1.05) {LM activations\\for same text\\middle layer $X$};

  \node[stage, text width=2.85cm] (measure) at (3.85,0.85) {\textbf{Measure alignment}\\fit ridge map\\$X\hat W \to Y$};
  \node[signal, text width=3.05cm] (score) at (3.85,-1.05) {alignment score\\unique $R^2$\\beyond nuisance};

  \node[stage, text width=2.95cm] (use) at (8.65,0.85) {\textbf{Use alignment}\\distill teacher\\to student};
  \node[loss, text width=2.85cm] (brainloss) at (8.65,-1.05) {$\mathcal{L}_{\mathrm{KD}} + \lambda\Lbrain$\\student budget};

  \draw[arrow] ([yshift=7pt]stim.east) to[out=0,in=180] (human.west);
  \draw[arrow] ([yshift=-7pt]stim.east) to[out=0,in=180] (lm.west);
  \draw[arrow] (human.east) to[out=0,in=180,looseness=1.15] (measure.west);
  \draw[arrow] (lm.east) to[out=0,in=195] ([yshift=-6pt]measure.west);
  \draw[arrow] (measure) -- (score);
  \draw[arrow] (measure.east) -- (use.west);
  \draw[arrow] (use) -- (brainloss);
  \draw[softarrow] (score.east) -- (brainloss.west);

  \node[draw=black!35, rounded corners=3pt, fit=(stim)(human)(lm), inner sep=7pt, label={[font=\footnotesize]above:input}] {};
  \node[draw=black!35, rounded corners=3pt, fit=(measure)(score), inner sep=7pt, label={[font=\footnotesize]above:measurement}] {};
  \node[draw=black!35, rounded corners=3pt, fit=(use)(brainloss), inner sep=7pt, label={[font=\footnotesize]above:intervention}] {};
\end{tikzpicture}
}
  \caption{The measurement instrument and its inversion. A ridge encoding model maps language-model activations to brain responses; held-out accuracy is reported as unique variance over the nuisance floor. The thesis asks whether a corresponding brain-alignment term can guide a distilled student.}
  \label{fig:pipeline}
\end{figure}

\textbf{Problem.}
We focus on model compression.
Knowledge distillation trains a small \emph{student} model to imitate a larger \emph{teacher}, reshaping the student's internal representations while it learns.
That is where the alignment question becomes sharp: if brain alignment lives in middle-layer structure, can a brain-alignment objective shape the student while it is being distilled \autocite{oota2023,oota2026}?
We ask whether such a student preserves or improves alignment at a matched compute budget, and whether any gain matters beyond the alignment score itself.
This case matters because compression is how small, deployable models are built, and recent brain-tuning and compression work has left this training use of brain alignment open \autocite{moussa2025,bilgin2026,merlin2026,oota2026}.

\textbf{Hypothesis.}
The hypothesis is that brain alignment can act, at most, as a weak prior over text-consistent representations.
The reason is a severe mismatch in scale.
A language model learns from on the order of a trillion tokens of text.
A language-fMRI dataset offers a few thousand noisy recordings, and a low noise ceiling caps the variance any model can explain in them \autocite{tuckute2024,lebel2023,lage2019}.
Against that imbalance, the brain term is too small to carry new task structure by itself.
Its plausible role is to bias the student among solutions already compatible with the text objective, the way a weak regularizer does \autocite{xu2017,cover2006,pirlot2022,freteault2025}.
This gives a falsifiable prediction.
A brain prior should help most where the text signal is weakest, in small, heavily compressed, or data-starved models, and should wash out wherever text data is abundant.
Section~\ref{sec:methods} makes the argument precise.

Recent studies have already used fMRI as an optimization signal for speech and text models \autocite{moussa2025,moussa2025b,negi2025,bilgin2026}.
The case is strengthened by causal evidence: text models pushed to be brain-\emph{misaligned} at matched perplexity perform worse across a broad downstream suite \autocite{merlin2026}.
Those studies live on the grow-and-optimize side of the field: they fine-tune the model, add adapters, or attach a brain-prediction head.
Compression work has taken the complementary path.
It shrinks a language model first, through quantization or pruning, and measures the alignment score afterward \autocite{oota2026}.
Alignment-guided distillation combines the two: the model is shrunk, and the alignment score enters the training objective.
That is the open case in Figure~\ref{fig:cells-diagram}, and it is exactly the regime where the weak-prior hypothesis says a brain signal has its best chance to matter.

\begin{figure}[t]
  \centering
  \resizebox{0.82\textwidth}{!}{\begin{tikzpicture}[
  font=\small,
  >=Latex,
  cell/.style={draw=black!55, rounded corners=3pt, fill=gray!6, minimum width=4.15cm, minimum height=1.42cm, text width=3.75cm, align=center, inner sep=5pt},
  thesis/.style={draw=orange!70!black, rounded corners=3pt, fill=orange!13, minimum width=4.15cm, minimum height=1.42cm, text width=3.75cm, align=center, inner sep=5pt, line width=0.8pt},
  head/.style={font=\small\bfseries, align=center},
  rowhead/.style={font=\small\bfseries, align=center, text width=2.4cm},
  note/.style={font=\footnotesize, align=center, text=black!70},
]
  \node[head] (c1) at (0,2.15) {Alignment measured\\after training};
  \node[head] (c2) at (4.45,2.15) {Alignment optimized\\during training};
  \node[rowhead] (r1) at (-3.4,0.85) {Model not\\shrunk};
  \node[rowhead] (r2) at (-3.4,-1.05) {Model\\shrunk};

  \node[cell] (gm) at (0,0.85) {encoding score\\\footnotesize frozen model};
  \node[cell] (go) at (4.45,0.85) {brain-guided tuning\\\footnotesize fine-tune / adapter / head};
  \node[cell] (sm) at (0,-1.05) {post-compression score\\\footnotesize quantization / pruning};
  \node[thesis] (so) at (4.45,-1.05) {alignment-guided\\distillation\\[-1pt]\footnotesize open in LM compression};
\end{tikzpicture}
}
  \caption{The field on two axes; three cells are occupied and the fourth, where this work sits, is open in language-model compression. Examples and citations are given in Table~\ref{tab:cells}.}
  \label{fig:cells-diagram}
\end{figure}

Brain-alignment scores are easy to inflate.
Shuffled splits leak through the temporal autocorrelation of fMRI, and low-level features such as sentence length carry much of an apparent score, so an untrained network can appear to explain brain data on its own \autocite{hadidi2026,jia2026}.
We therefore hold every number in this thesis to a fixed standard, set before measuring anything.
The contributions below summarize that standard, and Section~\ref{sec:methods} defines it \autocite{oota2023,merlin2024,oota2024,hadidi2026,jia2026}.

% Contributions list (beat 5) — Round 3 writes the text; scope = through Q2, methodology spine.
\paragraph{Contributions.}\leavevmode\par\smallskip
This thesis makes four contributions.
\begin{itemize}[noitemsep,topsep=2pt]
  \item It formulates alignment-guided distillation and an information-budget prediction for when a weak neural prior could matter.
  \item It establishes an artifact-resistant measurement protocol and a conditional within-subject result that survives whole-story splits, a capacity-matched nuisance floor, and an untrained-network control \evd{E006}.
  \item It shows that an averaged-target gain does not generalize to a per-individual effect, and tests that conclusion across adapter-capacity, loss-form, and voxelwise-mechanism variants \evd{E008}\evd{E011}\evd{E013}.
  \item It bounds practical payoff and synthetic-proxy transfer: no OOD-perplexity benefit is detected when the alignment manipulation is unreliable, a gaze-privileged route fails its preconditions, and a learnable dense neural proxy does not improve the common real-brain endpoint \evd{E009}\evd{E016}\evd{E024}.
\end{itemize}

The program asks five questions in sequence, each only if the previous one holds.
Is the signal real, or an artifact?
Does plain distillation preserve it?
Does an alignment objective move it beyond what perplexity already implies?
Does any gain survive per individual at matched perplexity?
And does it improve a downstream task?

The resulting arc is conditional rather than universal.
The signal survives the recorded controls within UTS03 \evd{E006}; plain distillation leaves alignment headroom but is quality-entangled \evd{E003}; the tested objective does not show a reliable fold-level gain \evd{E004}; its averaged-target positive does not generalize across participants \evd{E008}; and neither the tested practical routes nor the dense synthetic proxy establishes real-brain benefit \evd{E009}\evd{E016}\evd{E024}.
